\section{Proof of the Schur-field approximation}\label{app:schur-proof}

\begin{proof}[Proof of Lemma~\ref{lem:master-concentration}]
Fix a support $J$ and work on the range of $C_{JJ}$.  Every scalar
functional needed to net the support-specific terms in
\eqref{eq:master-sparse-norm} is a centered product of two linear forms
standardized by their population variances.  It therefore has uniformly
bounded $\psi_1$ norm.  A $1/4$-net in each support-specific whitened
sphere, followed by a union over $|J|\le2s$, has logarithmic size
$O(m)$.  Bernstein's inequality and the usual net approximation of
Euclidean and operator norms give \eqref{eq:master-concentration-tail},
including its linear term, for every $t\ge0$.  A matched Gaussian scalar functional
has the same variance as the corresponding original functional and is
sub-Gaussian, so the argument applies to every deterministic hybrid
pattern.

Integrating this Bernstein tail gives
\eqref{eq:master-concentration-moments} for each fixed $q\le6$ when
$\delta\le\delta_0$.  Taking the control-block component gives
\eqref{eq:master-relative-rip}.
\end{proof}
\begin{lemma}[Residual-field row bounds]
\label{lem:residual-field-row-bounds}
Use the notation and assumptions of
Theorem~\ref{thm:master-schur-field}, and put
\[
 L_S=C_{SS}^{-1}\Xi_S,\qquad
 w_S=U-L_S^\top X_S,\qquad
 R_S=\E(w_Sw_S^\top).
\]
Then
\begin{equation}\label{eq:residual-field-derivative-bound}
 \sup_{|S|\le s}\sup_{|H|_{U,s}\le1}
 |D\Phi_S(\Gamma_U)[H]|
 \lesssim\kappa_\varphi
 \left\{\lambda_s+\sup_{|S|\le s}|c_S-c_0|\right\}.
\end{equation}
If $Y=VV^\top-\Gamma_U$, $Y^G$ is a centered Gaussian matrix with the
covariance of $Y$, and $H$ is independent of both $Y$ and $Y^G$, then
\begin{align}
 \sup_{|S|\le s}
 \norm{D\Phi_S(\Gamma_U)[Y]+D^2\Phi_S(\Gamma_U)[H,Y]}_{\psi_1\mid H}
 &\lesssim\kappa_\varphi\left\{
 \lambda_s+\sup_{|S|\le s}|c_S-c_0|+|H|_{U,s}\right\},
 \label{eq:residual-field-mixed-row-bound}\\
 \sup_{|S|\le s}\left\{
 \norm{D^2\Phi_S(\Gamma_U)[Y,Y]}_{L_q}
 +\norm{D^2\Phi_S(\Gamma_U)[Y^G,Y^G]}_{L_q}\right\}
 &\lesssim\kappa_\varphi q(s+q),\qquad q\ge2.
 \label{eq:residual-field-diagonal-bound}
\end{align}
The mixed-row bound also holds with $Y^G$ in place of $Y$, and
\begin{equation}\label{eq:residual-field-diagonal-mean}
 \E D^2\Phi_S(\Gamma_U)[Y,Y]
 =\E D^2\Phi_S(\Gamma_U)[Y^G,Y^G].
\end{equation}
If $\varphi$ ignores an entry of its matrix argument, the corresponding
terms may be omitted from $|H|_{U,s}$.
\end{lemma}

\begin{proof}
Fix $S$ and write
\[
 H_{Sw}=H_{SU}-H_{SS}L_S,
 \qquad
 H_{ww}=H_{UU}-H_{US}L_S-L_S^\top H_{SU}
 +L_S^\top H_{SS}L_S.
\]
Differentiation on the compressed range gives
\begin{equation}\label{eq:residual-schur-derivatives}
\begin{split}
 D\mathcal R_S^U(\Gamma_U)[H]&=H_{ww},\\
 D^2\mathcal R_S^U(\Gamma_U)[H,Q]
 &=-H_{Sw}^\top C_{SS}^{-1}Q_{Sw}
 -Q_{Sw}^\top C_{SS}^{-1}H_{Sw}.
\end{split}
\end{equation}
For a centered row moment,
\begin{equation}\label{eq:row-residual-blocks}
 Y_{ww}=w_Sw_S^\top-R_S,
 \qquad Y_{Sw}=X_Sw_S^\top.
\end{equation}

Let $F_S=\nabla\varphi(R_S)$ and represent the Hessian by
$D^2\varphi(R_S)[P,Q]=\langle\mathcal H_S[P],Q\rangle$.  For one
unanchored support, the relevant row functional is
\begin{equation*}
\begin{split}
 &\langle F_S+\mathcal H_S[H_{ww}],Y_{ww}\rangle\\
 &-2\langle C_{SS}^{-1/2}H_{Sw}F_S,
 C_{SS}^{-1/2}Y_{Sw}\rangle_{\mathrm F}\\
 ={}&w_S^\top\{F_S+\mathcal H_S[H_{ww}]\}w_S
 -\operatorname{tr}[\{F_S+\mathcal H_S[H_{ww}]\}R_S]\\
 &-2X_S^\top C_{SS}^{-1}H_{Sw}F_Sw_S.
\end{split}
\end{equation*}
It is a centered quadratic form in $(U,X_S)$ whose coefficient has rank
at most six.

For the base derivative, put $K_S=(I_2,-L_S^\top)^\top$.  Its coefficient
in $(U,X_S)$ is
\[
 c_SK_SF_SK_S^\top-c_0K_\emptyset F_\emptyset K_\emptyset^\top.
\]
After block standardization by $\operatorname{diag}(D_U,C_{SS}^{1/2})$,
the $UU$, $UX$, and $XX$ blocks have operator norms bounded respectively
by
\[
 C\kappa_\varphi\{\lambda_s^2+|c_S-c_0|\},
 \qquad C\kappa_\varphi\lambda_s,
 \qquad C\kappa_\varphi\lambda_s^2.
\]
Indeed,
$R_S-R_\emptyset=-\Xi_S^\top C_{SS}^{-1}\Xi_S$ and
$C_{SS}^{1/2}L_SD_U^{-1}=C_{SS}^{-1/2}\Xi_SD_U^{-1}$.
Positive semidefiniteness of the joint covariance gives
$\lambda_s\le C$.  The standardized nuclear norm of the coefficient is
therefore bounded by
$C\kappa_\varphi\{\lambda_s+|c_S-c_0|\}$.  The same calculation and
\[
 \norm{C_{SS}^{-1/2}H_{Sw}D_U^{-1}}_{\mathrm F}
 +\norm{D_U^{-1}H_{ww}D_U^{-1}}_{\mathrm{op}}
 \le C|H|_{U,s}
\]
bound the standardized nuclear norm of the mixed coefficient by
$C\kappa_\varphi|H|_{U,s}$.  This proves
\eqref{eq:residual-field-derivative-bound}.  Relative sub-Gaussianity and
the spectral decomposition of a fixed-rank coefficient give the required
quadratic-form bound.  To make the standardization explicit, set
$\mathsf D_S=\operatorname{diag}(D_U,C_{SS}^{1/2})$ and
$\bar\Sigma_S=\mathsf D_S^{-1}\E[(U,X_S)(U,X_S)^\top]\mathsf D_S^{-1}$.
Positive semidefiniteness and the assumed target-block bound imply
$\norm{\bar\Sigma_S}_{\mathrm{op}}\le C$.  Hence, for every symmetric
coefficient $A$,
\[
 \norm{\bar\Sigma_S^{1/2}(\mathsf D_SA\mathsf D_S)
       \bar\Sigma_S^{1/2}}_*
 \le C\norm{\mathsf D_SA\mathsf D_S}_*.
\]
Applying the relative sub-Gaussian bound after this standardization, with
$z=(U^\top,X_S^\top)^\top$ and $\Sigma_z=\E(zz^\top)$, yields
\[
 \norm{z^\top Az-\E(z^\top Az)}_{\psi_1}
 \le CK^2\norm{\Sigma_z^{1/2}A\Sigma_z^{1/2}}_*,
\]
which proves \eqref{eq:residual-field-mixed-row-bound}.  The corresponding
functional of $Y^G$ is Gaussian with the same variance and obeys the same
bound.

For the diagonal term, the chain rule gives the exact identity
\begin{equation}\label{eq:residual-diagonal-formula}
\begin{split}
 D^2(\varphi\circ\mathcal R_S^U)(\Gamma_U)[Y,Y]
 ={}&D^2\varphi(R_S)[Y_{ww},Y_{ww}]\\
 &-2D\varphi(R_S)
 [Y_{Sw}^\top C_{SS}^{-1}Y_{Sw}].
\end{split}
\end{equation}
For an original row,
\[
 Y_{Sw}^\top C_{SS}^{-1}Y_{Sw}
 =(X_S^\top C_{SS}^{-1}X_S)w_Sw_S^\top.
\]
Relative sub-Gaussianity gives
\[
 \norm{X_S^\top C_{SS}^{-1}X_S}_{L_{2q}}\lesssim s+q,
 \qquad
 \norm{D_U^{-1}w_Sw_S^\top D_U^{-1}}_{L_{2q}}\lesssim q.
\]
Together with \eqref{eq:master-standardized-derivatives}, this proves the
first term in \eqref{eq:residual-field-diagonal-bound}.

For the matched Gaussian row, put
\[
 \mathcal Z_S=C_{SS}^{-1/2}(Y^G)_{Sw}D_U^{-1}.
\]
If $\norm A_{\mathrm F}=1$, covariance matching and a singular-value
decomposition of the $|S|\times2$ matrix $A$ give
\[
 \Var\langle A,\mathcal Z_S\rangle
 =\Var\langle A,C_{SS}^{-1/2}X_Sw_S^\top D_U^{-1}\rangle
 \le C.
\]
Thus the covariance operator of $\operatorname{vec}(\mathcal Z_S)$ is
bounded and its trace is $O(s)$.  Gaussian quadratic-form moments yield
$\norm{\mathcal Z_S^\top\mathcal Z_S}_{L_q}\lesssim s+q$.
The fixed-dimensional block
$D_U^{-1}(Y^G)_{ww}D_U^{-1}$ has $L_{2q}$ norm $O(\sqrt q)$, so
\eqref{eq:residual-diagonal-formula} proves the Gaussian part of
\eqref{eq:residual-field-diagonal-bound}.  Finally,
\eqref{eq:residual-field-diagonal-mean} follows because the bilinear form
$D^2\Phi_S(\Gamma_U)$ is deterministic and $Y,Y^G$ have the same
covariance.
\end{proof}

\begin{proof}[Proof of Theorem~\ref{thm:master-schur-field}]
Let $Y_i=V_iV_i^\top-\Gamma_U$, and let $Y_i^G$ be independent centered
Gaussian matrices with the covariance of $Y_i$.
Lemma~\ref{lem:master-concentration} supplies the required fixed moments
and tails uniformly over deterministic hybrid rows
$\widetilde Y_i\in\{Y_i,Y_i^G,0\}$.
The same differentiation used in
\eqref{eq:residual-schur-derivatives}, the derivative neighborhood, and
\eqref{eq:master-standardized-derivatives} give, for
$|H|_{U,s}\le c$,
\begin{equation}\label{eq:generic-higher-derivatives}
 \sup_{|S|\le s}\sup_{0\le t\le1}
 \left|D^k\Phi_S(\Gamma_U+tH)[H,\ldots,H]\right|
 \le C\kappa_\varphi|H|_{U,s}^k,
 \qquad k=2,3,
\end{equation}
and hence
\begin{equation}\label{eq:generic-taylor-remainder}
 \sup_{|S|\le s}\left|
 \Phi_S(\Gamma_U+H)-\Phi_S(\Gamma_U)-D\Phi_S(\Gamma_U)[H]
 -\frac12D^2\Phi_S(\Gamma_U)[H,H]
 \right|\le C\kappa_\varphi|H|_{U,s}^3.
\end{equation}
The derivative bound in
\eqref{eq:residual-field-derivative-bound},
\eqref{eq:generic-higher-derivatives}, and
\eqref{eq:master-concentration-tail} prove
\eqref{eq:master-uniform-deviation}--\eqref{eq:master-second-order}.
For $H$ independent of a centered
row moment $Y$ and $q\ge2$, Lemma~\ref{lem:residual-field-row-bounds}
gives
\begin{align}
 \sup_{|S|\le s}
 \norm{D\Phi_S(\Gamma_U)[Y]+D^2\Phi_S(\Gamma_U)[H,Y]}_{\psi_1\mid H}
 &\lesssim\kappa_\varphi
 \{\lambda_s+\Delta_c+|H|_{U,s}\},
 \label{eq:generic-linear-row-bound}\\
 \sup_{|S|\le s}
 \norm{D^2\Phi_S(\Gamma_U)[Y,Y]}_{L_q}
 &\lesssim\kappa_\varphi q(s+q).
 \label{eq:generic-diagonal-row-bound}
\end{align}
Both bounds hold for $Y^G$, and
\begin{equation}\label{eq:generic-diagonal-mean-matching}
        \E D^2\Phi_S(\Gamma_U)[Y,Y]
        =\E D^2\Phi_S(\Gamma_U)[Y^G,Y^G],
\end{equation}
by \eqref{eq:residual-field-diagonal-mean}.

We now compare the two Taylor fields.  For either row law, put
\[
        \mathsf{Diag}_{n,S}
        =\frac1{2n^2}\sum_{i=1}^n
        \{D^2\Phi_S(\Gamma_U)[Y_i,Y_i]
        -\E D^2\Phi_S(\Gamma_U)[Y,Y]\}.
\]
The quadratic Taylor term decomposes exactly as
\begin{equation}\label{eq:generic-taylor-diagonal-split}
\begin{split}
 \frac12D^2\Phi_S(\Gamma_U)\left[\frac1n\sum_iY_i,
                         \frac1n\sum_iY_i\right]
 ={}&\frac1{n^2}\sum_{i<j}D^2\Phi_S(\Gamma_U)[Y_i,Y_j]\\
 &+\frac1{2n}\E D^2\Phi_S(\Gamma_U)[Y,Y]+\mathsf{Diag}_{n,S}.
\end{split}
\end{equation}
For centered independent $Z_i$ with the moment growth in
\eqref{eq:generic-diagonal-row-bound}, the independent-sum inequality
\[
 \norm{\textstyle\sum_{i=1}^n Z_i}_{L_\nu}
 \lesssim
 \sup_{2\le q\le\nu}
 \frac\nu q\left(\frac n\nu\right)^{1/q}\norm{Z_1}_{L_q}
 \lesssim \kappa_\varphi
 \{s\sqrt{n\nu}+\nu(s+\nu)\}
\]
follows by symmetrization and the standard moment bound for sums.  A
union over at most $e^m$ supports then gives, for $\nu=m+u\lesssim n$,
\begin{equation}\label{eq:generic-diagonal-field-concentration}
 \Prob\left[
 \max_{|S|=s}|\mathsf{Diag}_{n,S}|
 >C\kappa_\varphi
 \left\{\frac{s\sqrt\nu}{n^{3/2}}
 +\frac{\nu(s+\nu)}{n^2}\right\}
 \right]\le e^{-c\nu}.
\end{equation}
Since there are at most $e^m$ supports, the moment bound preceding
\eqref{eq:generic-diagonal-field-concentration}, applied with
$\nu_0=m\vee2$, also gives directly
\[
 \left\|\max_{|S|=s}|\mathsf{Diag}_{n,S}|\right\|_{L_{\nu_0}}
 \lesssim \kappa_\varphi
 \left\{\frac{s\sqrt{\nu_0}}{n^{3/2}}
 +\frac{\nu_0(s+\nu_0)}{n^2}\right\}.
\]
Here the factor from the maximum is at most
$e^{m/\nu_0}\le e$.  Because
$a_{n,\Phi}\ge\kappa_\varphi m/n$, $s\le m$, and
$\nu_0\asymp m$, division by $a_{n,\Phi}$ bounds the right side by
\[
 C\left\{\frac{s}{\sqrt{nm}}+\frac{s+m}{n}\right\}
 \le C\delta.
\]
Thus the centered diagonal field is $O(\delta)$ after normalization in
$L_1$, and hence also for the truncated loss.  Its retained mean agrees
for the two row laws by \eqref{eq:generic-diagonal-mean-matching}.

It remains to replace the off-diagonal field.  At the $i$th step of a
row-by-row replacement, let
\[
        H_i=\frac1n\sum_{j\ne i}\widetilde Y_j.
\]
Conditional on $H_i$, inserting a row $Y$ changes the normalized
support value by the linear functional
\[
        \Delta_{i,S}(Y)
        =\frac{D\Phi_S(\Gamma_U)[Y]
        +D^2\Phi_S(\Gamma_U)[H_i,Y]}{na_{n,\Phi}}.
\]
The original and Gaussian increments have the same conditional mean and
covariance.  Equation~\eqref{eq:generic-linear-row-bound} and the
definition of $a_{n,\Phi}$ give
\begin{equation}\label{eq:generic-off-diagonal-row-bound}
        \sup_{|S|=s}\norm{\Delta_{i,S}(Y)}_{\psi_1\mid H_i}
        +\sup_{|S|=s}\norm{\Delta_{i,S}(Y^G)}_{\psi_1\mid H_i}
        \lesssim\frac{\Lambda_i}{\sqrt{mn}},
        \qquad
        \Lambda_i=1+\frac{|H_i|_{U,s}}{\delta}.
\end{equation}
The bound also holds for signed combinations with bounded total
variation.  Equations~\eqref{eq:master-concentration-moments} and
\eqref{eq:master-concentration-tail} imply bounded fixed moments of
$\Lambda_i$.  Applying the latter with
$t\asymp n^{2/3}m^{1/3}=o(n)$ also gives
\[
        \sum_{i=1}^n\Prob\{\Lambda_i>c(n/m)^{1/3}\}
        =o\{(m/n)^{1/6}\}.
\]

Approximate the maximum by log-sum-exp with inverse temperature
$\lambda m$; the approximation error is at most $1/\lambda$.  A
second-order Taylor expansion in
$\Delta_i$ cancels conditionally through order two.  Let $\pi_S$ be the
Gibbs weights before the insertion and $\pi_S(\theta)$ those after
inserting $\theta\Delta_i$.  Convexity of the log-partition function gives,
for $k\le3$,
\[
 \prod_{a=1}^k\pi_{S_a}(\theta)
 \le \prod_{a=1}^k\pi_{S_a}
 \exp\!\left[\theta\lambda m\left\{
 \sum_{a=1}^k\Delta_{i,S_a}
 -k\sum_T\pi_T\Delta_{i,T}\right\}\right].
\]
The expression in braces is a signed combination with conditional
$\psi_1$ norm at most $C\Lambda_i/\sqrt{mn}$.  On
$\Lambda_i\le c(n/m)^{1/3}$, the choice
$\lambda=(n/m)^{1/6}$ makes its tilted exponential moment bounded, since
$\lambda\sqrt{m/n}\,\Lambda_i\le c$.  Conditional H\"older inequalities
and \eqref{eq:generic-off-diagonal-row-bound} then bound the one-row
third-order remainder by
\[
 C\{1+\lambda m+(\lambda m)^2\}
 \frac{\Lambda_i^3}{(mn)^{3/2}}.
\]
After summing the rows and adding the smoothing error, the bound is
\[
        C\left\{
        \frac1\lambda
        +\frac1{m^{3/2}\sqrt n}
        +\frac{\lambda}{\sqrt{mn}}
        +\lambda^2\sqrt{\frac mn}\right\}.
\]
Taking $\lambda=(n/m)^{1/6}$ proves an
$O\{(m/n)^{1/6}\}$ comparison between the off-diagonal fields.

Finally, take $t$ to be a sufficiently small fixed multiple of $n$ in
\eqref{eq:master-concentration-tail}, so that
$H=\hat\Gamma_U-\Gamma_U$ lies in the Taylor neighborhood outside an
event of probability $Ce^{-cn}$.  On this event,
\eqref{eq:generic-taylor-remainder} and the fixed-moment form of sparse
concentration give
\[
 \E\left[1\wedge
 \frac{\kappa_\varphi|H|_{U,s}^3}{a_{n,\Phi}}\right]
 \le C\delta,
\]
because $a_{n,\Phi}\ge\kappa_\varphi\delta^2$; bounded test functions
absorb the exceptional event.  Together with
\eqref{eq:generic-diagonal-field-concentration}, this proves
\eqref{eq:master-matched}.  Subtracting one fixed support changes each
bound by at most a constant factor and proves \eqref{eq:master-anchored}.
\end{proof}

\section{Product approximation}\label{app:product-proofs}

\begin{lemma}[Score-product remainder]\label{lem:score-product-error}
Suppose the assumptions of Theorem~\ref{thm:chaos-degeneration} hold.
Let $R_n^G$ and $R_n^E$ denote the two suprema on the
left side of \eqref{eq:score-reduction}.  Then
\begin{equation}\label{eq:score-product-error-moments}
 \E\left(1\wedge\frac{R_n^G}{a_n^{\mathrm{score}}}\right)
 +\E\left(1\wedge\frac{R_n^E}{a_n^{\mathrm{score}}}\right)
 \lesssim\delta.
\end{equation}
\end{lemma}

\begin{proof}
The polynomial remainder bound in the proof of
\eqref{eq:score-reduction}, with the cubic Taylor remainder added for the
Gaussian field, controls $R_n^G$ and $R_n^E$.  Applying the fixed moments
and full tail in Lemma~\ref{lem:master-concentration} to these bounds proves
\eqref{eq:score-product-error-moments}.
\end{proof}

\begin{proof}[Proof of
Theorem~\ref{thm:gaussian-top-s-decoupling}]
In standardized coordinates, set
\[
 F(z)=\max_{|S|=s}-z_{1S}^\top z_{2S},
 \qquad
 F_\tau(z)=\frac1\tau
 \log\sum_{|S|=s}\exp\{-\tau z_{1S}^\top z_{2S}\}.
\]
The smoothing error is at most $m/\tau$.  Set
$\Omega_t=(1-t)I_{2p}+t\Omega$.  Gaussian interpolation along
$\{\Omega_t:0\le t\le1\}$, followed by integration by parts,
expresses the derivative of $\E F_\tau$ as one half of the inner product
of $\Omega-I_{2p}$ with $\E\nabla^2F_\tau$.  Each supportwise quadratic
Hessian is supported on one set of size $s$, and the softmax covariance
of two supportwise gradients is supported on their union.  Sparse score
isotropy and \eqref{eq:score-mean-cap} therefore give
\[
 \left|\frac{d}{dt}\E F_\tau\right|
 \lesssim\rho_s^{\mathrm{score}}(s+\tau m).
\]
Taking $\tau=(\rho_s^{\mathrm{score}})^{-1/2}$ when the covariance error
is nonzero shows that the two unsmoothed expectations differ by $o(m)$.
The almost-everywhere gradient of $F$ is supported on a maximizing set,
where
\[
 \nabla F^\top\Omega\nabla F
 \le(1+\rho_s^{\mathrm{score}})\norm{\nabla F}^2.
\]
Gaussian Poincar\'e and the sparse score bound give variance $O(m)$
under both laws.  Independent placement of the two fields on one
probability space therefore gives a difference
$o_{\Prob}(m)$; restoring the marginal scales proves
\eqref{eq:gaussian-score-decoupling}.
\end{proof}

\section{Studentized score geometry}\label{app:studentized-geometry}

For $|S|\le s$, define
\begin{equation}\label{eq:studentized-directions}
\begin{gathered}
        u_{n,S}=\begin{cases}
        M_S\mathbf y/\norm{M_S\mathbf y}_2,
        &\norm{M_S\mathbf y}_2>0,\\
        0,&\norm{M_S\mathbf y}_2=0,
        \end{cases}
        \qquad
        \nu_{n,S}
        =\sqrt{\frac{n-|S|-1}{n-\operatorname{rank}(\mathbf X_S)}},\\
        u_{n,S}^{\Delta}
        =\nu_{n,S}u_{n,S}
        -\nu_{n,\emptyset}u_{n,\emptyset}.
\end{gathered}
\end{equation}

\begin{lemma}[Score-to-direction identity]
\label{lem:studentized-score-direction}
For supports $S,T$ with nonzero outcome residuals,
\begin{align}
        (I-M_S)\mathbf y&=\frac{\mathbf X\hat q_{n,S}}{\sqrt n}, \notag\\
        \norm{M_S\mathbf y}_2^2
        &=\norm{\mathbf y}_2^2-\norm{\hat q_{n,S}}_{\hat C}^2,
        \label{eq:score-projection-identity}\\
 \norm{u_{n,S}^{\Delta}-u_{n,T}^{\Delta}}_2^2
 &= (\nu_{n,S}-\nu_{n,T})^2 \notag\\
 &\quad+\frac{\nu_{n,S}\nu_{n,T}}
 {\norm{M_S\mathbf y}_2\norm{M_T\mathbf y}_2}
 \left[
 \norm{\hat q_{n,S}-\hat q_{n,T}}_{\hat C}^2
 -\{\norm{M_S\mathbf y}_2-\norm{M_T\mathbf y}_2\}^2
 \right].
 \label{eq:studentized-score-direction-bridge}
\end{align}
\end{lemma}

\begin{proof}
The Moore--Penrose identity gives
\[
        \frac{\mathbf X\hat q_{n,S}}{\sqrt n}
        =\mathbf X_S(\mathbf X_S^\top\mathbf X_S)^\dagger
        \mathbf X_S^\top\mathbf y=(I-M_S)\mathbf y.
\]
Pythagoras gives the residual-norm identity in
\eqref{eq:score-projection-identity}.  Moreover,
\[
        \norm{M_S\mathbf y-M_T\mathbf y}_2
        =\norm{\hat q_{n,S}-\hat q_{n,T}}_{\hat C}.
\]
For nonzero vectors $a,b$,
\[
 \left\|\frac a{\norm a}-\frac b{\norm b}\right\|^2
 =\frac{\norm{a-b}^2-(\norm a-\norm b)^2}
 {\norm a\,\norm b}.
\]
Apply this identity to $a=M_S\mathbf y$ and $b=M_T\mathbf y$, then use
\[
 \norm{\nu_{n,S}u_{n,S}-\nu_{n,T}u_{n,T}}_2^2
 =(\nu_{n,S}-\nu_{n,T})^2
 +\nu_{n,S}\nu_{n,T}\norm{u_{n,S}-u_{n,T}}_2^2.\qedhere
\]
\end{proof}

\begin{lemma}[Score packing implies direction packing]
\label{lem:score-direction-packing}
Suppose \eqref{eq:studentized-score-packing} holds and
$\delta\le\delta_0$ for a sufficiently small fixed constant.  Then
\begin{equation}\label{eq:studentized-direction-packing}
 \min_{\substack{S,T\in\mathcal P\\S\ne T}}
 \norm{u_{n,S}^{\Delta}-u_{n,T}^{\Delta}}_2
 \ge c\frac{d_n}{\norm{\mathbf y}_2},
 \qquad
 \max_{S\in\mathcal P}\norm{u_{n,S}^{\Delta}}_2
 \le C\frac{d_n}{\norm{\mathbf y}_2}.
\end{equation}
\end{lemma}

\begin{proof}
For $|S|\le s$, small $\delta$ gives
\[
 0<c\le\nu_{n,S}\le C,
 \qquad
 |\nu_{n,S}-\nu_{n,T}|\le C\frac{s+1}{n}.
\]
Equations~\eqref{eq:studentized-score-direction-bridge} and
the angular condition in \eqref{eq:studentized-score-packing}, together with
$r_{n,S}r_{n,T}\le\norm{\mathbf y}_2^2$, give the separation in
\eqref{eq:studentized-direction-packing}.  Taking $T=\emptyset$ in
\eqref{eq:studentized-score-direction-bridge}, dropping the nonpositive
radial term in its bracket, and using
the residual floor in \eqref{eq:studentized-score-packing} gives
\[
 \norm{u_{n,S}^{\Delta}}_2
 \le C\left\{\frac{s+1}{n}
 +\frac{\norm{\hat q_{n,S}}_{\hat C}}{\norm{\mathbf y}_2}\right\}.
\]
The score-radius and rank-scale bounds in
\eqref{eq:studentized-score-packing} prove the radius bound.
\end{proof}

\begin{lemma}[Conditional studentized-direction reduction]
\label{lem:conditional-studentized-direction}
Let $\mathcal F=\sigma(\mathbf X,\mathbf y)$, and suppose that the entries
of $\mathbf x_0\in\R^n$ are independent of $\mathcal F$, independent
across observations, centered, have variance $\sigma_0^2$, and satisfy
$\norm{x_0^{(i)}}_{\psi_2}\le K\sigma_0$.  Put
\begin{equation}\label{eq:conditional-treatment-scale}
        \hat\sigma_0^2=\frac{\mathbf x_0^\top\mathbf x_0}{n},
        \qquad
        \Delta^x_{n,S}
        =\frac{\mathbf x_0^\top M_S\mathbf x_0/
        \{n-\operatorname{rank}(\mathbf X_S)\}}
        {\mathbf x_0^\top\mathbf x_0/n}-1.
\end{equation}
If $\delta\le\delta_0$, then, outside an event of probability $Ce^{-cm}$,
\begin{equation}\label{eq:conditional-studentized-direction-reduction}
 \hat\sigma_0\asymp\sigma_0,
 \qquad
 \max_{|S|=s}\left|
 \hat t_{0,S}-\hat t_{0,\emptyset}
 -\frac{\mathbf x_0^\top u_{n,S}^{\Delta}}{\hat\sigma_0}
 \right|
 \le C\sqrt m\,\delta^2.
\end{equation}
\end{lemma}

\begin{proof}
For every support with
$\mathbf x_0^\top M_S\mathbf x_0>0$ and
$\norm{M_S\mathbf y}_2>0$, the partial-correlation identity gives exactly
\begin{equation}\label{eq:exact-studentized-direction}
        \hat t_{0,S}
        =\nu_{n,S}\frac{\mathbf x_0^\top u_{n,S}}{\hat\sigma_0}
        (1+\Delta^x_{n,S})^{-1/2}.
\end{equation}
If $M_S\mathbf y=0$, both sides are zero under the stated convention.

Conditional quadratic and linear-form concentration, followed by a union
over the supports, gives
\begin{equation}\label{eq:conditional-direction-event}
        \hat\sigma_0\asymp\sigma_0,
        \qquad
        \max_{|S|=s}|\Delta^x_{n,S}|\le C\delta^2,
        \qquad
        \max_{|S|=s}|\mathbf x_0^\top u_{n,S}|
        \le C\sigma_0\sqrt m
\end{equation}
outside an event of probability $Ce^{-cm}$.  To see the quadratic bound
explicitly, let $P_S=I-M_S$.  Since $P_S$ has rank at most $s$, the same
conditional union bound gives
\[
 \max_{|S|=s}\mathbf x_0^\top P_S\mathbf x_0
 \le C\sigma_0^2m,
 \qquad
 \mathbf x_0^\top\mathbf x_0\ge cn\sigma_0^2.
\]
One has
\[
 1+\Delta^x_{n,S}
 =\frac{n}{n-\operatorname{rank}(\mathbf X_S)}\left(1-
 \frac{\mathbf x_0^\top P_S\mathbf x_0}
 {\mathbf x_0^\top\mathbf x_0}\right),
\]
which proves the middle assertion in
\eqref{eq:conditional-direction-event}.  The final assertion is the
corresponding conditional union bound over the unit vectors $u_{n,S}$.

Since $\Delta^x_{n,\emptyset}=0$,
\eqref{eq:exact-studentized-direction} and
$|(1+z)^{-1/2}-1|\le C|z|$ for $|z|\le1/2$ give
\[
 \max_{|S|=s}\left|
 \hat t_{0,S}-\hat t_{0,\emptyset}
 -\frac{\mathbf x_0^\top u_{n,S}^{\Delta}}{\hat\sigma_0}
 \right|
 \le C\sqrt m\,\delta^2
\]
on the event in \eqref{eq:conditional-direction-event}.
\end{proof}
